\documentclass[10pt]{extarticle}
\usepackage[a4paper,margin=1.5cm]{geometry}
\usepackage[T1]{fontenc}
\usepackage{amsmath,amssymb}
\usepackage{graphicx}
\usepackage[export]{adjustbox}
\usepackage{multicol}
\usepackage{enumitem}
\usepackage{booktabs,array,tabularx}
\setlength{\columnsep}{0.7cm}
\setlength{\parindent}{0pt}
\setlength{\parskip}{3pt}
\pagestyle{empty}
\makeatletter
% figure registry (figdefs.tex is generated by the extract step)
\newcommand{\qfigdef}[4]{\@namedef{qf@file@#1}{#2}\@namedef{qf@w@#1}{#3}\@namedef{qf@wide@#1}{#4}}
\newcommand{\qf@check}[1]{\@ifundefined{qf@file@#1}{\PackageError{score2}{Unknown figure id #1}{}}{}}
\newcommand{\qf@img}[4][]{\includegraphics[width=#2\dimexpr\@nameuse{qf@w@#4}pt\relax,max width=#3,max height=0.45\textheight,#1]{\@nameuse{qf@file@#4}}}
% Full-width questions: finalize wraps every question that contains a wide figure in \qfullbegin...\qfullend,
% so the WHOLE question (stem, figures, options) is set across the page in reading order, then the two columns
% resume. The question's head - its number, stem and first figure (\qfullsplit marks the end of the head) -
% is kept on one page, so a number is never left at a page foot with its figure overleaf; the rest may break
% across pages normally, so no big gaps are left. Without a \qfullsplit, the whole question is kept together
% if it is under 45% of a page. Never write these yourself.
\newif\ifqf@full
\newif\ifqf@boxed
\newsavebox{\qf@box}
\newcommand{\qf@open}{\global\qf@fulltrue\par\global\qf@boxedtrue
  \setbox\qf@box\vbox\bgroup\hsize\textwidth\linewidth\textwidth}
\newcommand{\qf@emit}[1]{\par\egroup\global\qf@boxedfalse
  \ifdim\dimexpr\ht\qf@box+\dp\qf@box\relax<#1\textheight\noindent\usebox{\qf@box}\par\else\unvbox\qf@box\fi}
\newcommand{\qfullsplit}{\ifqf@boxed\qf@emit{0.85}\fi}
\newcommand{\qf@close}{\ifqf@boxed\qf@emit{0.45}\fi\global\qf@fullfalse}
\newcommand{\qfullbegin}{\end{multicols}\qf@open}
\newcommand{\qfullend}{\qf@close\begin{multicols}{2}}
\newcommand{\qfullstart}{\qf@open}
\newcommand{\qfulldone}{\qf@close}
% \qfig[scale]{id}: block figure, centred, directly where it is written (size set by the tool).
% Fallback if a wide figure is not inside a full-width question: break out of the columns in place - no floats.
\newcommand{\qfig}[2][1]{\qf@check{#2}%
  \ifqf@full\qf@blk{#1}{#2}\else
  \ifnum\@nameuse{qf@wide@#2}=1 \expandafter\@firstoftwo\else\expandafter\@secondoftwo\fi
  {\end{multicols}\begin{center}\qf@img{#1}{\textwidth}{#2}\end{center}\begin{multicols}{2}}%
  {\qf@blk{#1}{#2}}\fi}
\newcommand{\qf@blk}[2]{\par\smallskip{\centering\qf@img{#1}{\linewidth}{#2}\par}\smallskip}
% \qfigi[scale]{id}: inline figure (inside option lists / tables / minipages); never breaks columns.
\newcommand{\qfigi}[2][1]{\qf@check{#2}\qf@img[valign=c]{#1}{\linewidth}{#2}}
\makeatother
% structure macros - lint reads these to check completeness and figure placement
\newcommand{\papersec}[1]{\par\bigskip{\centering\Large\bfseries PAPER-#1\par}\medskip}
\newcommand{\partsec}[2]{\par\medskip\noindent{\large\bfseries PART-#1: #2}\par\smallskip}
\newcommand{\sect}[1]{\par\smallskip\noindent\textbf{#1}\par}
\newcommand{\qn}[1]{\par\smallskip\noindent\textbf{#1.}~}
\newcommand{\paraq}[2]{\par\smallskip\noindent\textbf{Paragraph for Questions #1 and #2}\par}
\newenvironment{qoptions}{\begin{itemize}[leftmargin=1.8em,itemsep=1pt,topsep=2pt,label={}]}{\end{itemize}}
\newcommand{\opt}[2]{\item[(#1)] #2}
% figure registry (generated by score2_tools.py finalize)
\qfigdef{p1_phy_s1_q01_1}{Score Advanced - Test 09 (29 Apr 2023)-img/p1_phy_s1_q01_1.png}{175.8}{0}
\qfigdef{p1_phy_s1_q02_B}{Score Advanced - Test 09 (29 Apr 2023)-img/p1_phy_s1_q02_B.png}{124.4}{0}
\qfigdef{p1_phy_s1_q04_1}{Score Advanced - Test 09 (29 Apr 2023)-img/p1_phy_s1_q04_1.png}{111.5}{0}
\qfigdef{p1_phy_s2_q05_1}{Score Advanced - Test 09 (29 Apr 2023)-img/p1_phy_s2_q05_1.png}{108.3}{0}
\qfigdef{p1_phy_s2_q06_1}{Score Advanced - Test 09 (29 Apr 2023)-img/p1_phy_s2_q06_1.png}{82.0}{0}
\qfigdef{p1_phy_s2_q07_1}{Score Advanced - Test 09 (29 Apr 2023)-img/p1_phy_s2_q07_1.png}{108.3}{0}
\qfigdef{p1_phy_s2_q08_1}{Score Advanced - Test 09 (29 Apr 2023)-img/p1_phy_s2_q08_1.png}{92.6}{0}
\qfigdef{p1_phy_s2_q09_1}{Score Advanced - Test 09 (29 Apr 2023)-img/p1_phy_s2_q09_1.png}{108.3}{0}
\qfigdef{p1_phy_s2_q10_1}{Score Advanced - Test 09 (29 Apr 2023)-img/p1_phy_s2_q10_1.png}{108.3}{0}
\qfigdef{p1_phy_s2_q11_1}{Score Advanced - Test 09 (29 Apr 2023)-img/p1_phy_s2_q11_1.png}{133.2}{0}
\qfigdef{p1_phy_s2_q12_1}{Score Advanced - Test 09 (29 Apr 2023)-img/p1_phy_s2_q12_1.png}{162.3}{0}
\qfigdef{p1_phy_s3_q02_1}{Score Advanced - Test 09 (29 Apr 2023)-img/p1_phy_s3_q02_1.png}{106.6}{0}
\qfigdef{p1_phy_s3_q03_1}{Score Advanced - Test 09 (29 Apr 2023)-img/p1_phy_s3_q03_1.png}{216.2}{0}
\qfigdef{p1_phy_s3_q04_1}{Score Advanced - Test 09 (29 Apr 2023)-img/p1_phy_s3_q04_1.png}{245.2}{0}
\qfigdef{p1_phy_s3_q05_1}{Score Advanced - Test 09 (29 Apr 2023)-img/p1_phy_s3_q05_1.png}{39.7}{0}
\qfigdef{p1_phy_s3_q06_1}{Score Advanced - Test 09 (29 Apr 2023)-img/p1_phy_s3_q06_1.png}{174.9}{0}
\qfigdef{p1_chem_s1_q03_1}{Score Advanced - Test 09 (29 Apr 2023)-img/p1_chem_s1_q03_1.png}{81.8}{0}
\qfigdef{p1_chem_s1_q04_A}{Score Advanced - Test 09 (29 Apr 2023)-img/p1_chem_s1_q04_A.png}{107.5}{0}
\qfigdef{p1_chem_s1_q04_B}{Score Advanced - Test 09 (29 Apr 2023)-img/p1_chem_s1_q04_B.png}{107.8}{0}
\qfigdef{p1_chem_s1_q04_D}{Score Advanced - Test 09 (29 Apr 2023)-img/p1_chem_s1_q04_D.png}{179.9}{0}
\qfigdef{p1_chem_s2_q05_A}{Score Advanced - Test 09 (29 Apr 2023)-img/p1_chem_s2_q05_A.png}{132.5}{0}
\qfigdef{p1_chem_s2_q05_B}{Score Advanced - Test 09 (29 Apr 2023)-img/p1_chem_s2_q05_B.png}{132.7}{0}
\qfigdef{p1_chem_s2_q07_1}{Score Advanced - Test 09 (29 Apr 2023)-img/p1_chem_s2_q07_1.png}{240.1}{0}
\qfigdef{p1_chem_s2_q10_C}{Score Advanced - Test 09 (29 Apr 2023)-img/p1_chem_s2_q10_C.png}{137.6}{0}
\qfigdef{p1_chem_s2_q10_D}{Score Advanced - Test 09 (29 Apr 2023)-img/p1_chem_s2_q10_D.png}{97.7}{0}
\qfigdef{p1_chem_s2_q11_C}{Score Advanced - Test 09 (29 Apr 2023)-img/p1_chem_s2_q11_C.png}{141.3}{0}
\qfigdef{p1_chem_s3_q05_1}{Score Advanced - Test 09 (29 Apr 2023)-img/p1_chem_s3_q05_1.png}{217.5}{0}
\qfigdef{p1_chem_s3_q05_3}{Score Advanced - Test 09 (29 Apr 2023)-img/p1_chem_s3_q05_3.png}{92.3}{0}
\qfigdef{p1_chem_s3_q06_1}{Score Advanced - Test 09 (29 Apr 2023)-img/p1_chem_s3_q06_1.png}{69.2}{0}
\qfigdef{p1_chem_s3_q06_2}{Score Advanced - Test 09 (29 Apr 2023)-img/p1_chem_s3_q06_2.png}{316.1}{1}
\qfigdef{p2_phy_s1_q01_1}{Score Advanced - Test 09 (29 Apr 2023)-img/p2_phy_s1_q01_1.png}{92.0}{0}
\qfigdef{p2_phy_s1_q02_1}{Score Advanced - Test 09 (29 Apr 2023)-img/p2_phy_s1_q02_1.png}{145.4}{0}
\qfigdef{p2_phy_s1_q03_1}{Score Advanced - Test 09 (29 Apr 2023)-img/p2_phy_s1_q03_1.png}{154.8}{0}
\qfigdef{p2_phy_s1_q04_1}{Score Advanced - Test 09 (29 Apr 2023)-img/p2_phy_s1_q04_1.png}{192.8}{0}
\qfigdef{p2_phy_s1_q05_1}{Score Advanced - Test 09 (29 Apr 2023)-img/p2_phy_s1_q05_1.png}{97.9}{0}
\qfigdef{p2_phy_s1_q06_1}{Score Advanced - Test 09 (29 Apr 2023)-img/p2_phy_s1_q06_1.png}{147.9}{0}
\qfigdef{p2_phy_s1_q07_1}{Score Advanced - Test 09 (29 Apr 2023)-img/p2_phy_s1_q07_1.png}{102.6}{0}
\qfigdef{p2_phy_s1_q08_1}{Score Advanced - Test 09 (29 Apr 2023)-img/p2_phy_s1_q08_1.png}{179.0}{0}
\qfigdef{p2_phy_s2_q09_1}{Score Advanced - Test 09 (29 Apr 2023)-img/p2_phy_s2_q09_1.png}{77.7}{0}
\qfigdef{p2_phy_s2_q10_1}{Score Advanced - Test 09 (29 Apr 2023)-img/p2_phy_s2_q10_1.png}{82.0}{0}
\qfigdef{p2_phy_s2_para11_12_1}{Score Advanced - Test 09 (29 Apr 2023)-img/p2_phy_s2_para11_12_1.png}{94.2}{0}
\qfigdef{p2_phy_s3_q01_1}{Score Advanced - Test 09 (29 Apr 2023)-img/p2_phy_s3_q01_1.png}{68.0}{0}
\qfigdef{p2_phy_s3_q02_1}{Score Advanced - Test 09 (29 Apr 2023)-img/p2_phy_s3_q02_1.png}{74.7}{0}
\qfigdef{p2_phy_s3_q03_1}{Score Advanced - Test 09 (29 Apr 2023)-img/p2_phy_s3_q03_1.png}{161.1}{0}
\qfigdef{p2_phy_s3_q04_1}{Score Advanced - Test 09 (29 Apr 2023)-img/p2_phy_s3_q04_1.png}{145.0}{0}
\qfigdef{p2_phy_s3_q05_1}{Score Advanced - Test 09 (29 Apr 2023)-img/p2_phy_s3_q05_1.png}{130.6}{0}
\qfigdef{p2_phy_s3_q06_1}{Score Advanced - Test 09 (29 Apr 2023)-img/p2_phy_s3_q06_1.png}{102.4}{0}
\qfigdef{p2_chem_s1_q04_1}{Score Advanced - Test 09 (29 Apr 2023)-img/p2_chem_s1_q04_1.png}{115.1}{0}
\qfigdef{p2_chem_s1_q04_2}{Score Advanced - Test 09 (29 Apr 2023)-img/p2_chem_s1_q04_2.png}{108.0}{0}
\qfigdef{p2_chem_s2_para9_10_1}{Score Advanced - Test 09 (29 Apr 2023)-img/p2_chem_s2_para9_10_1.png}{245.2}{0}
\qfigdef{p2_math_s2_q09_1}{Score Advanced - Test 09 (29 Apr 2023)-img/p2_math_s2_q09_1.png}{145.2}{0}
\qfigdef{p2_math_s2_q10_1}{Score Advanced - Test 09 (29 Apr 2023)-img/p2_math_s2_q10_1.png}{154.6}{0}
\raggedcolumns
\makeatletter\@beginparpenalty=10000\makeatother
\newcolumntype{L}[1]{>{\raggedright\arraybackslash}p{#1}}
\begin{document}
\begin{center}{\Large\bfseries Score Advanced --- Test 09 (29 Apr 2023)}\end{center}
\begin{multicols}{2}
\raggedcolumns
\papersec{1}
\partsec{1}{PHYSICS}
\sect{SECTION-I (i)}
\qn{1} The network shown consists of an ideal battery, an ideal voltmeter V, an ideal ammeter A, several resistors of different values and a rheostat. If resistance of the rheostat is increased by a finite amount, how will the readings of the voltmeter and ammeter change?\par
\qfig{p1_phy_s1_q01_1}
\begin{qoptions}
\opt{A}{Both will decrease}
\opt{B}{Voltmeter reading will decrease and ammeter reading will increase.}
\opt{C}{Voltmeter reading will increase and ammeter reading will decrease.}
\opt{D}{Voltmeter reading will not change but ammeter reading will decrease.}
\end{qoptions}
\qn{2} A solid hemispherical bowl of radius R and mass m rests on the frictionless plane (Assumed to be X-Z plane) as shown in the figure. A force F = mg is applied on the tip of the hemisphere as shown in the figure. The x-component of acceleration of point having coordinate of (R, 0.625 R, 0) will be\par
\qfig{p1_phy_s1_q02_B}
\begin{qoptions}
\opt{A}{2g}
\opt{B}{$\sqrt{2}$ g}
\opt{C}{$\dfrac{\text{g}\sqrt{3}}{2}$}
\opt{D}{g}
\end{qoptions}
\qn{3} A body starts moving on x-axis from rest from x = $-$2 m. Acceleration $\vec{\text{a}}$ of the particle varies according to the equation\par
$\vec{\text{a}} = -(\pi\text{ rad/sec})^2\,(\vec{\text{x}})$ for 0 $\leqslant$ t $\leqslant$ $\dfrac{1}{2}$sec\par
$\vec{\text{a}} = \left(5\pi\text{ m/s}^2\right)\left(\dfrac{\vec{\text{x}}}{|\vec{\text{x}}|}\right)$ for t $>$ $\dfrac{1}{2}$sec.\par
where $\vec{\text{x}}$ is the position vector of the particle at any instant t.\par
Then velocity of the particle at x = $\dfrac{\pi}{2}$m is
\begin{qoptions}
\opt{A}{2 $\pi$ m/s}
\opt{B}{3 $\pi$ m/s}
\opt{C}{5 $\pi$ m/s}
\opt{D}{none of these}
\end{qoptions}
\qn{4} In the given figure two concentric cylindrical region in which time varying magnetic field is present as shown. From the centre to radius R magnetic field is perpendicular into the plane varying as $\dfrac{\text{dB}}{\text{dt}}$ = 2k$_0$ and in a region from R to 2R magnetic field is perpendicular out of the plane varying as $\dfrac{\text{dB}}{\text{dt}}$ = 4k$_0$. Find the induced emf across an arc AB of radius 3R.\par
\qfig{p1_phy_s1_q04_1}
\begin{qoptions}
\opt{A}{6R$^2$k$\theta$}
\opt{B}{5R$^2$k$\theta$}
\opt{C}{7R$^2$k$\theta$}
\opt{D}{None of these}
\end{qoptions}
\sect{SECTION-I (ii)}
\qn{5} A plank of mass M=4kg is symmetrically placed over two identical cylinders of radius R and identical mass M as shown in the following figure. A vertical spring of natural length R and spring constant k = 700 $\times$ $\pi^2$ N/m is connecting the center of the plank with the ground. All surfaces are rough and there is no slippage anywhere. The plank is displaced horizontally by a very small length A from the mean position and released. Choose the correct option(s).\par
\qfig{p1_phy_s2_q05_1}
\begin{qoptions}
\opt{A}{Motion of the plank is simple harmonic.}
\opt{B}{Motion of the system is periodic.}
\opt{C}{Time period of oscillation is 0.7s}
\opt{D}{Time period of oscillation is 0.2s}
\end{qoptions}
\qn{6} In the following system a massless L shape body is constraint to move freely along the horizontal surface. A cylinder of mass M = 4kg and radius R = 2m is attached as shown in the figure. Coefficient of kinetic friction of the cylinder with the vertical surface of the L shaped body is $\mu$ = 0.2. The L shaped body is pulled left by applying 200N force as shown. Choose the correct option(s).\par
\qfig{p1_phy_s2_q06_1}
\begin{qoptions}
\opt{A}{Acceleration of the L shaped body is 80/3 m/s$^2$}
\opt{B}{Tension in the upper string is 110 N}
\opt{C}{Frictional force at the vertical surface of L shaped body is 40 N}
\opt{D}{Angular acceleration of the cylinder is 45/2 rad/s$^2$}
\end{qoptions}
\qn{7} Two tubes of comparable radii and lengths x and y respectively are joined as shown and maintained at different uniform temperatures T$_1$ = $-$153\textdegree C and T$_2$ = 207\textdegree C respectively (ignore any end correction). Choose the correct option(s):\par
\qfig{p1_phy_s2_q07_1}
\begin{qoptions}
\opt{A}{For join of tube to be pressure Node y/x may be 2}
\opt{B}{For join of tube to be pressure Node y/x may be 4}
\opt{C}{For join of tube to be pressure antinode y/x may be 2}
\opt{D}{For join of tube to be pressure antinode y/x may be 6}
\end{qoptions}
\qn{8} Interference pattern experiment is done on the surface of water by using a slit in an opaque plate AB and screen is placed at 50 cm from slit. Slit is 1 cm above water surface and light of wavelength $\lambda$ = 6000 \AA{} is falling over it. The plate is oscillating according to equation x = A sin $\omega$t along the BO, where A = 1 cm, $\omega$ = 100 rad/s. Point O is just above the water surface and another point P is 5 cm above the water surface. Plate AB is just above the water surface.\par
\qfig{p1_phy_s2_q08_1}
\begin{qoptions}
\opt{A}{Maximum speed of fringe at point P is 10 cm/s}
\opt{B}{Maximum speed of fringe at point P is 5 cm/s}
\opt{C}{Intensity at point O is minimum}
\opt{D}{Intensity at point O is maximum}
\end{qoptions}
\qn{9} An ideal gas having degree of freedom f is kept inside a thermally insulated vessel. Vessel having two large pistons each with charge of +q and $-$q. Assume charges are uniformly distributed. At an instant charge at pistons is increased by $\beta$ times with help of an external agent. Then which of the following statements are correct ? No friction between wall and piston.\par
\qfig{p1_phy_s2_q09_1}
\begin{qoptions}
\opt{A}{Work done by the agent in the process is ($\beta^2-1$)nRT$_1$ where T$_1$ is initial temp.}
\opt{B}{Work done by the agent in the process is ($\beta^2-1$)$\dfrac{\text{q}^2\text{L}_1}{\text{A}\in_0}$ where L$_1$ is initial separation between pistons}
\opt{C}{The ratio of final to initial temperature of gas in equilibrium is $\dfrac{2\beta^2+\text{f}}{2+\text{f}}$}
\opt{D}{The ratio of final to initial temperature of gas in equilibrium is $\dfrac{\beta^2+\text{f}}{2+\text{f}}$}
\end{qoptions}
\qn{10} A solid paraboloid of base radius R$_0$ and height h$_0$ is having uniform volume mass density of $\rho_\text{s}$ is inverted and placed just above the surface as shown, when paraboloid is released from rest it has been observed that when it becomes completely submerged in the liquid the paraboloid comes to the rest again. Assume liquid level is not changing as paraboloid enters in liquid. Liquid having volume mass density is $\rho_\text{L}$. (g is uniform)\par
\qfig{p1_phy_s2_q10_1}
\begin{qoptions}
\opt{A}{The phenomenon is true for all paraboloids irrespective of their dimensions as long as density ratio is maintained.}
\opt{B}{The ratio of $\dfrac{\rho_\text{L}}{\rho_\text{s}}$ is 2.}
\opt{C}{The ratio of $\dfrac{\rho_\text{L}}{\rho_\text{s}}$ is 3.}
\opt{D}{The phenomenon is not true for all paraboloids irrespective of their dimensions as long as density ratio is maintained.}
\end{qoptions}
\qn{11} The figure shows equipotential contours for two charges located along the X-axis. Two points A and B have been marked on 1.2 volt and 2.7 volt contours respectively. Choose the correct statement(s).\par
\qfig{p1_phy_s2_q11_1}
\begin{qoptions}
\opt{A}{The charges are in the ratio 1:3.}
\opt{B}{Both the charges are positive.}
\opt{C}{Potential energy of a negative charge increases when moved from point A to point B}
\opt{D}{On X-axis, the field is zero some where between x = $-$1 and x = +1.}
\end{qoptions}
\qn{12} For an experiment of photoelectric effect, a graph between 1/$\lambda$ and stopping potential (V) of three metals A, B and C whose work functions are $\phi_\text{A}$, $\phi_\text{B}$ \& $\phi_\text{C}$ respectively, is plotted as shown in the figure. Which of the following statement(s) is/are correct? [Here $\lambda$ is the wavelength of incident light]\par
\qfig{p1_phy_s2_q12_1}
\begin{qoptions}
\opt{A}{$\theta_\text{A}<\theta_\text{B}<\theta_\text{C}$}
\opt{B}{$\theta_\text{A}=\theta_\text{B}=\theta_\text{C}$}
\opt{C}{$\phi_\text{A}<\phi_\text{B}<\phi_\text{C}$}
\opt{D}{$\phi_\text{A}=\phi_\text{B}=\phi_\text{C}$}
\end{qoptions}
\sect{SECTION-II}
\qn{1} Bottom face of a horizontal dielectric slab is coated with a conducting paint and at centre of the top face is placed a small mercury drop, which does not wet the plate. Thickness of the slab is d and relative permittivity is $\varepsilon_\text{r}$. Now the switch is closed and the battery voltage is increased so gradually that the mercury drop spreads on the dielectric slab at negligible rate. At what battery voltage will the mercury drop cover the whole of the upper face of the dielectric? Surface tension of mercury is $\sigma$ and initial surface area of the drop is negligible as compared to surface area of mercury when spread. If your answer is V = $\sqrt{\dfrac{\text{x}\sigma\text{d}}{\varepsilon_\text{r}\varepsilon_0}}$ write value of x.
\qn{2} The circuit consists of a capacitor, inductance, and two resistors, see figure. The voltage on both resistors is 10 V, and the voltage between the leads A and B is also 10 V. Find the applied voltage U$_0$ (in volt).\par
\qfig{p1_phy_s3_q02_1}
\qn{3} A large block of mass m$_\text{b}$ = 5kg is located on a horizontal frictionless surface. A second block of mass m$_\text{t}$ = 2kg is located on top of the first block. The coefficient of friction (both static and kinetic) between the two blocks is given by 0.5. All surfaces are horizontal and all motion is effectively one dimensional. A spring with spring constant k = 500 N/m is connected to the top block only. The spring obeys Hooke's Law equally in both extension and compression. Assume that the top block never falls off the bottom block. You may assume that the bottom block is very, very long. The top block is moved a distance A away from the equilibrium position and then released from rest. Depending on the value of A, the motion can be divided into two types: motion that experiences no frictional energy losses and motion that does. Find the value A$_\text{c}$ (in cm) that divides the two motion types.\par
\qfig{p1_phy_s3_q03_1}
\qn{4} In the given figure two co-axial cylinders are placed as shown. Inner cylinder of radius r contains one mole of an ideal gas enclosed by a piston which can move in inner cylinder without friction and the space between the cylinder of radius 2r and radius r contains liquid. Heat is supplied to the gas kept in inner cylinder at a constant rate of 100 R J/cycle by some external agent. The cyclic process for the gas is represented by the indicator diagram as shown. Now if heat lost by the gas equal to heat gain by the liquid and radiation is negligible, find the rate of increasing of temperature of liquid per cycle (mass of liquid is 1 kg and specific heat capacity = 10 R in SI unit, temperature at points 1, 2 and 3 in cycle are 300 K, 400 K and 600 K respectively, R is universal gas constant).\par
\qfig{p1_phy_s3_q04_1}
\qn{5} A glass tube nearly filled with mercury is attached to the bottom of a light iron rod of length 1 m (see figure). How high (in m) must the mercury be in the glass tube for the centre of mass of this pendulum not to rise or fall with change in temperature?\par
($\alpha_\text{iron}$ = 12 $\times$ 10$^{-6}$ K$^{-1}$, $\gamma_\text{mercury}$ = 18.2 $\times$ 10$^{-5}$ K$^{-1}$ and $\alpha_\text{glass}$ = 9 $\times$ 10$^{-6}$ K$^{-1}$)\par
\qfig{p1_phy_s3_q05_1}
\qn{6} A large plane mirror with its bottom on the floor is tilted at an angle 30\textdegree{} to the vertical. A boy whose eyes are at height $\sqrt{3}$ m above the floor is standing in front of the mirror. At what maximum distance (in m) from mirror should the boy be to see his full image in mirror?\par
\qfig{p1_phy_s3_q06_1}
\partsec{2}{CHEMISTRY}
\sect{SECTION-I (i)}
\qn{1} The coloured salt P was warmed with NaOH solution and colourless gas Q was evolved and a coloured solution R was produced. The colour of salt P was different from that of solution R. On adding BaCl$_2$ to R, a coloured PPT was obtained. Then the correct formula for salt P is
\begin{qoptions}
\opt{A}{Na$_2$Cr$_2$O$_7$}
\opt{B}{Mg(NH$_4$)PO$_4$}
\opt{C}{(NH$_4$)$_2$Cr$_2$O$_7$}
\opt{D}{PH$_4$I}
\end{qoptions}
\qn{2} A XF$_4$ molecule have $\mu$ = 0, which additional informational is required to confirm the geometry of molecule.
\begin{qoptions}
\opt{A}{All the X$-$F bond length are equal}
\opt{B}{Molecule has same F$-$X$-$F bond Angle between any two adjacent F}
\opt{C}{Number of lone pair of electrons on central atom $\leq$ 2}
\opt{D}{Planarity of molecule}
\end{qoptions}
\qn{3} For the equilibrium 2NH$_3$(g) $\rightleftharpoons$ N$_2$(g) + 3H$_2$(g) of reaction starts with pure NH$_3$. Density d(gm/Lit.) of equilibrium mixture varies with degree of dissociation ($\alpha$) of NH$_3$ of 237\textdegree C in such a manner that $\dfrac{1}{\text{d}}$ v/s ``$\alpha$'' is a linear plot shown below, what is the total pressure at equilibrium.\par
\qfig{p1_chem_s1_q03_1}
\begin{qoptions}
\opt{A}{5.110 atm}
\opt{B}{3.862 atm}
\opt{C}{2.463 atm}
\opt{D}{1.56 atm}
\end{qoptions}
\qn{4} Select the CORRECT option.
\begin{qoptions}
\opt{A}{\qfigi{p1_chem_s1_q04_A} can be differentiated by Victor Mayer's Test}
\opt{B}{\qfigi{p1_chem_s1_q04_B} can be differentiated by Beilstein test}
\opt{C}{Monomer of polyethene is Vinyl chloride}
\opt{D}{\qfigi{p1_chem_s1_q04_D} can be distinguished by Seliwanoff's test}
\end{qoptions}
\sect{SECTION-I (ii)}
\qn{5} K$_2$[PtCl$_4$] + 2Kgly $\longrightarrow$ X + 4KCl\par
[Pt(NH$_3$)$_4$](NO$_3$)$_2$ + 2Kgly $\longrightarrow$ Y + 2KNO$_3$ + 4NH$_3$\par
Here Kgly is potassium glycinate and `X' is soluble in polar solvent and `Y' is soluble in non-polar solvent. Then
\begin{qoptions}
\opt{A}{\qfigi{p1_chem_s2_q05_A}}
\opt{B}{\qfigi{p1_chem_s2_q05_B}}
\opt{C}{Both X and Y are optically inactive}
\opt{D}{Both X and Y are in trans form}
\end{qoptions}
\qn{6} Cu on reaction with dil. HNO$_3$ liberates a colourless gas which of the following is FALSE for the gas?
\begin{qoptions}
\opt{A}{Has 11 valance e$^-$, where extra electron occupy an antibonding 2p orbital}
\opt{B}{An acid anhydride}
\opt{C}{is isoelectronic with superoxide}
\opt{D}{Atomicity is 2}
\end{qoptions}
\qn{7} Consider the process carried out for an ideal gas.\par
\qfig{p1_chem_s2_q07_1}
One by one, stoppers S$_1$, S$_2$, S$_3$ ...... are removed slowly and the gas expands to V$_2$. L with final pressure P$_2$ and final temperature T$_2$. Which of the following is/are correct for the above process.
\begin{qoptions}
\opt{A}{T$_1$ = T$_2$}
\opt{B}{P$_1$ = P$_2$}
\opt{C}{U$_1$ = U$_2$}
\opt{D}{w = 0, q = 0}
\end{qoptions}
\qn{8} An ideal binary liquid solution of A and B has a vapour pressure of 900 mm of Hg. It is distilled at constant temperature till $\dfrac{2}{3}$ of the original amount is distilled. If the mole fraction of `A' in the residue and mole fraction of B in the condensates 0.3 and 0.4 respectively and vapour pressure of residue is 860 mm of Hg, then which of the following is/are correct options :
\begin{qoptions}
\opt{A}{The initial mixture taken should be an equimolar mixture of A and B}
\opt{B}{The first vapors formed will have more moles of A as compared to moles of B.}
\opt{C}{P$^\circ_\text{A}$ = 1000 torr}
\opt{D}{The vapor pressure above the condensate is 920 torr at same temperature}
\end{qoptions}
\qn{9} NH$_3$ concentration is a saturated solution of AgBr(s) is maintained at 0.2M. The correct information(s) related with the solution is/are :\par
K$_\text{sp}$ of Ag Br = 2 $\times$ 10$^{-10}$\par
K$_\text{f}$ of Ag(NH$_3$)$_2^+$ = 8 $\times$ 10$^6$\par
$\lambda^\text{o}_{\text{(Ag}^+)}$ = 6 $\times$ 10$^{-3}$ Sm$^2$ mol$^{-1}$\par
$\lambda^\text{o}_{\text{(Br}^-)}$ = 8 $\times$ 10$^{-3}$ Sm$^2$ mol$^{-1}$
\begin{qoptions}
\opt{A}{Conductivity of Br$^-$ = 6.4 $\times$ 10$^{-2}$ S m$^{-1}$}
\opt{B}{Conductivity of Ag$^+$ = 1.5 $\times$ 10$^{-7}$ S m$^{-1}$}
\opt{C}{[Ag$^+$] = 8 $\times$ 10$^{-3}$M}
\opt{D}{[Br$^-$] = 2 $\times$ 10$^{-3}$M}
\end{qoptions}
\qn{10} Which of the given reactions follows E1$_\text{cb}$ mechanism?
\begin{qoptions}
\opt{A}{CF$_3-$CHCl$_2$ $\xrightarrow{\text{RO}^\ominus}$}
\opt{B}{Ph$-$CH$_2-$CH$_2-$F $\xrightarrow{\text{C}_2\text{H}_5\text{O}^\ominus}$}
\opt{C}{\qfigi{p1_chem_s2_q10_C}}
\opt{D}{\qfigi{p1_chem_s2_q10_D}}
\end{qoptions}
\qn{11} Correctly matched product in the given reactions is/are :
\begin{qoptions}
\opt{A}{CH$\equiv$CH $\xrightarrow[\text{(i)NH}_4\text{Cl, (ii)HCl(1eq), (iii)Polymerisation}]{\text{Cu}_2\text{Cl}_2}$ Neoprene}
\opt{B}{CH$_3-\overset{\text{(R)}}{\text{CH}}-\text{CO}_2\text{H}$ with Cl on the chiral carbon, $\xrightarrow[\text{(ii) NaNO}_2+\text{HCl}]{\text{(i) NH}_3}$ CH$_3\overset{\text{(R)}}{\text{CH}}-\text{CO}_2\text{H}$ with OH on the chiral carbon}
\opt{C}{\qfigi{p1_chem_s2_q11_C}}
\opt{D}{CH$_3-$C$\equiv$C$-$CH$_3$ $\xrightarrow[\substack{\text{\scriptsize(ii) EtI}\\\text{\scriptsize(iii) H}_2\text{\scriptsize(Pd--BaSO}_4\text{\scriptsize)}\\\text{\scriptsize(iv) OsO}_4\text{\scriptsize, NaHSO}_3}]{\text{\scriptsize(i) NaNH}_2\text{\scriptsize{} in paraffin solvent}}$ Meso - hexane - 3,4diol}
\end{qoptions}
\qn{12} Acyl halides can be reduced to aldehydes by treatment with `X'.\par
$\text{R}-\overset{\text{O}}{\overset{\|}{\text{C}}}-\text{Cl} \xrightarrow{\text{`X'}} \text{R}-\overset{\text{O}}{\overset{\|}{\text{C}}}-\text{H}$\par
`X' can be
\begin{qoptions}
\opt{A}{Lithium tri-t-butoxyaluminium Hydride in diglyme at $-$78\textdegree C}
\opt{B}{1 Eq. Nickel as the catalyst during hydrogenolysis.}
\opt{C}{Hydrogenolysis with Pd-BaSO$_4$}
\opt{D}{BH$_3$$-$THF/ACOH}
\end{qoptions}
\sect{SECTION-II}
\qn{1} The conductivity of 0.001M Na$_2$SO$_4$ solution is 2.6 $\times$ 10$^{-4}$ S cm$^{-1}$ and it increases to 7 $\times$ 10$^{-4}$ S cm$^{-1}$, when the solution is saturated with CaSO$_4$. The molar conductivity of Na$^+$ and Ca$^{2+}$ are 50 and 120 S cm$^2$ mol$^{-1}$, respectively. Neglecting conductivity of water, the solubility product of CaSO$_4$ is Y $\times$ 10$^{-\text{x}}$ M$^2$. Then y is.
\qn{2} Metal M of radius 50 nm is crystallise in FCC type and made up of unit cell such that face of unit cells are aligned with face of cubic crystal. If the total number of metal atoms of M at all faces of cubical crystal is 6 $\times$ 10$^{30}$, then the area of one face of cubical crystal is A $\times$ 10$^{16}$ m$^2$, find A.
\qn{3} How many of the following have been correctly matched with their property.\par
(a) ClO$_2$ $-$ Contains odd electron but does not undergo dimerisation.\par
(b) XeF$_8^{-2}$ $-$ Does not exist.\par
(c) Cl$_2$O $-$ Bond angle greater than 109\textdegree 28'.\par
(d) O$_3^-$ $-$ Paramagnetic\par
(e) POCl$_3$ $-$ ClPCl bond angle less than 109\textdegree 28'.\par
(f) N$_2$H$_4$ $-$ N$-$N bond length greater than N$_2$F$_4$\par
(g) H$_2$S$_2$O$_8$ $-$ Contains peroxy linkage.\par
(h) XeF$_4$ $-$ d$_{\text{x}^2-\text{y}^2}$ orbital is involved in hybridisation of central atom.
\qn{4} Compound A on treatment with NaIO$_3$ in HCl; gives a violet coloured substance (B). The violet colour disappears on treatment with (C), producing A and D. Compound `C' is formed by reaction between sulphur and sodium sulphite. Find the maximum number of bonds between identical atoms having equal bond length in compound D.
\qn{5} Consider the following reactions :\par
\qfigi{p1_chem_s3_q05_1}\par
(b) $\text{CH}_3-\text{CH}_2-\overset{\text{O}}{\overset{\|}{\text{C}}}-\text{H} \xrightarrow[5\text{\textdegree C}]{\text{NaOH}} \text{Y}$ (Number of aldol product)\par
(c) $\text{CH}_3-\text{CH}_2-\overset{\text{O}}{\overset{\|}{\text{C}}}-\text{OCH}_3 \xrightarrow[\text{H}^+]{\text{CH}_3\text{O}^\ominus\text{K}} \text{Z}$ (Number of $\beta$-keto ester product)\par
\qfigi{p1_chem_s3_q05_3}\par
Find out the X + Y + Z + W?
\qfullbegin
\qn{6} How many of the following will produce Aromatic product?\par
(A) $\text{CH}_3-\text{C}\equiv\text{CH} \xrightarrow[873\text{K}]{\text{Red hot iron Fe tube}}$ \quad (B) $\text{C}_6\text{H}_{14} \xrightarrow[\text{Cr}_2\text{O}_3/\Delta]{\text{Al}_2\text{O}_3}$ \quad (C) \qfigi{p1_chem_s3_q06_1}\par
(D) Benzene $\xrightarrow[\text{Dark/Cold}]{\text{Cl}_2\text{,AlCl}_3\text{ excess}}$\par
\qfig{p1_chem_s3_q06_2}
\qfullend
\partsec{3}{MATHEMATICS}
\sect{SECTION-I (i)}
\qn{1} A box has 18 balls, 6 white balls numbered 1 to 6, 6 black balls numbered 1 to 6 and 6 red balls numbered 1 to 6. We have to divide all 18 balls into 6 groups (groups are not numbered) such that each group consists of one white, one black and one red ball and difference of numbers on any two balls of same group is at most one. Find no. of ways this division can be done.
\begin{qoptions}
\opt{A}{97}
\opt{B}{98}
\opt{C}{99}
\opt{D}{100}
\end{qoptions}
\qn{2} Consider two circles in a plane with radii 5 and 3 respectively. Distance between centres of these two circles is 10. If R$_1$ is radius of circle passing through 4 points of contact of direct common tangents of these two circles and R$_2$ is radius of circle passing through 4 points of contact of transverse common tangents then $\dfrac{\text{R}_1}{\text{R}_2}$ =?
\begin{qoptions}
\opt{A}{$\dfrac{\sqrt{5}+\sqrt{3}}{\sqrt{5}-\sqrt{3}}$}
\opt{B}{2}
\opt{C}{$\dfrac{5}{3}$}
\opt{D}{$\dfrac{25}{9}$}
\end{qoptions}
\qn{3} For x, y, z $\in$ [0, 2$\pi$]. The number of ordered triplets (x, y, z) satisfying 16($\sin^4$x) + 2($\sin^4$y) + 4($\sin^4$z) $-$ 16(sin x) (sin y) (sin z) + 2 = 0.
\begin{qoptions}
\opt{A}{20}
\opt{B}{18}
\opt{C}{16}
\opt{D}{14}
\end{qoptions}
\qn{4} Let f(x) be a linear function such that f(0) = $-$ 5 and f(f(0)) = $-$15. The number of values of m for which the solutions of the inequality f(x)f(m $-$ x) $\geq$ 0 form an interval of length 2, is.
\begin{qoptions}
\opt{A}{0}
\opt{B}{1}
\opt{C}{2}
\opt{D}{3}
\end{qoptions}
\sect{SECTION-I (ii)}
\qn{5} If A and B are two matrices of order 3 $\times$ 3 such that 2A + 3BB$^\text{T}$ = I and B$^{-1}$ = A$^\text{T}$, where A$^\text{T}$ represents transpose of matrix A also T$_\text{r}$ (A) represents trace of matrix A and det(A) represents determinant of matrix A, then
\begin{qoptions}
\opt{A}{A and B are both symmetric matrix}
\opt{B}{det(A$^{-1}$ $-$ 3B$^3$ + BA) = 27}
\opt{C}{Tr(A$^{-1}$ $-$ 3B$^3$ $-$ BA) = 3}
\opt{D}{A and B are commute of each other}
\end{qoptions}
\qn{6} Consider f(x) = \par\noindent\resizebox{\linewidth}{!}{$\begin{vmatrix}5&7&3\\2&9g^2(\text{x})&2\text{x}^2\\\text{x}&3g(\text{x})-\text{x}^2&1\end{vmatrix}-\begin{vmatrix}2&4&5\\2\text{x}^2&2&9g^2(\text{x})\\1&\text{x}&3g(\text{x})-\text{x}^2\end{vmatrix}+\begin{vmatrix}2&\text{x}&0\\9g^2(\text{x})&3g(\text{x})-\text{x}^2&-2\\2\text{x}^2&1&-1\end{vmatrix}$}\par
(where g(x) is polynomial function) such that m = $\displaystyle\int_0^1 f(\text{x})\text{dx}$ is minimum. Then
\begin{qoptions}
\opt{A}{m = 1/5}
\opt{B}{m = 2/5}
\opt{C}{Number of point of inflection of $\phi(\text{x}) = \displaystyle\int_0^\text{x}\left[\left(\int_0^\text{t} g(\text{y})\text{dy}-\left(\text{e}^\text{t}-\text{e}^{-\text{t}}\right)\right)\right]\text{dt}$ is 0}
\opt{D}{g(x) $\leqslant$ cos x $\forall$ x $\in$ [$-$1, 1]}
\end{qoptions}
\qn{7} n digit numbers are formed using the digit 1, 2, 3, 4 each of which begins with 1 such that no two consecutive digit are same. If a$_\text{n}$(n $\geqslant$ 2) be such n-digit numbers which end with 1 while b$_\text{n}$ be the ones which ends in a fixed number other than 1, then
\begin{qoptions}
\opt{A}{value of a$_7$ is 183}
\opt{B}{2a$_{11}$ = 2a$_{10}$ + 3a$_9$}
\opt{C}{value of a$_5$ is 21}
\opt{D}{a$_{11}$ = 2a$_{10}$ + 3a$_9$}
\end{qoptions}
\qn{8} Consider the planes\par
P$_1$ : x = y sin $\alpha$ + z sin $\beta$\par
P$_2$ : y = z sin $\theta$ + x sin $\alpha$\par
P$_3$ : z = x sin $\beta$ + y sin $\theta$\par
If P$_1$, P$_2$, P$_3$ intersect in a line $\dfrac{\text{x}}{\cos\theta}=\dfrac{\text{y}}{\cos\beta}=\dfrac{\text{z}}{\cos\alpha}$ then
\begin{qoptions}
\opt{A}{$\sin^2\alpha+\sin^2\beta+2\sin\theta\sin\alpha\sin\beta=\cos^2\theta$}
\opt{B}{$\sin^3\alpha+\sin^3\beta+\sin^3\theta=3\sin\alpha\sin\beta\sin\theta$}
\opt{C}{$\theta+\alpha+\beta=\dfrac{\pi}{2}$}
\opt{D}{$\theta+\alpha+\beta=\pi$}
\end{qoptions}
\qn{9} A chord of negative slope from P $\left(\sqrt{264},\,0\right)$ is drawn to ellipse x$^2$ + 4y$^2$ = 16 this chord intersects the ellipse at A and B, O is origin then
\begin{qoptions}
\opt{A}{maximum area of $\Delta$AOB is 4}
\opt{B}{If area of $\Delta$AOB is maximum then slope of AB is $\dfrac{-1}{4\sqrt{2}}$}
\opt{C}{maximum area of $\Delta$AOB is 8}
\opt{D}{If area of $\Delta$AOB is maximum then slope of AB is $\dfrac{-1}{8\sqrt{2}}$}
\end{qoptions}
\qn{10} For a $\Delta$ABC, vertex A $\equiv$ (4, 4), vertex B lies on L$_1$ : y $-$ 2 = 0 and vertex C lies on L$_2$ : x $-$ y + 4 = 0 if perimeter of $\Delta$ABC is minimum then which of following options is/are correct.
\begin{qoptions}
\opt{A}{B $\equiv$ (4, 2) and C $\equiv$ (0, 4)}
\opt{B}{B $\equiv$ (3, 2) and C $\equiv$ $\left(\dfrac{4}{3},\dfrac{16}{3}\right)$}
\opt{C}{Area of $\Delta$ABC is equal to 4 square units}
\opt{D}{Minimum Perimeter of $\Delta$ABC is 6 + 2$\sqrt{5}$}
\end{qoptions}
\qn{11} Define the collections \{E$_1$, E$_2$, E$_3$, .....\} of ellipse and \{R$_1$, R$_2$, R$_3$, .....\} of rectangles as follows :\par
E$_1$ : $\dfrac{\text{x}^2}{25}+\dfrac{\text{y}^2}{9}=1$;\par
R$_1$ : Rectangle of largest area with sides parallel to the axes, inscribed in E$_1$.\par
E$_\text{n}$ : Ellipse $\dfrac{\text{x}^2}{\text{a}_\text{n}^2}+\dfrac{\text{y}^2}{\text{b}_\text{n}^2}=1$ of largest area inscribed in R$_{\text{n}-1}$, n $>$ 1, then which of the following options is/are correct?
\begin{qoptions}
\opt{A}{The eccentricities of E$_{15}$ and E$_{16}$ are equal}
\opt{B}{The length of latus rectum of E$_7$ is $\dfrac{9}{20}$}
\opt{C}{$\displaystyle\sum_{\text{n}=1}^{\text{N}}$ (area of R$_\text{n}$) $>$ 60, for each positive integer N}
\opt{D}{The distance between focus of ellipse E$_{11}$ is $\dfrac{1}{4}$}
\end{qoptions}
\qn{12} For a non-zero complex number Z, let Re(Z) denotes the real part of Z, let S be the set of all complex numbers Z satisfying Z$^4$ $-$ $|$Z$|^4$=12i(Z)$^2$ where i = $\sqrt{-1}$. If Z$_1$, Z$_2$ $\in$ S and Re(Z$_1$) $>$ 0, Re(Z$_2$) $<$ 0 then which of the following options is/are correct?
\begin{qoptions}
\opt{A}{Minimum value of $|\text{Z}_1-\text{Z}_2|$ is $\sqrt{6}$}
\opt{B}{Minimum value of $|\text{Z}_1-\text{Z}_2|$ is 2$\sqrt{6}$}
\opt{C}{Area of triangle formed by any tangent of locus of Z with coordinate axes is 6}
\opt{D}{Area of triangle formed by any tangent of locus of Z with coordinate axes is 12}
\end{qoptions}
\sect{SECTION-II}
\qn{1} Let f(x) be a polynomial of degree three such that f(0) = 5, f(1) = 8 and x = 0 is a critical point of f(x) such that f(x) does not have a local extremum at x = 0. If $\displaystyle\int_{-1}^{1}\dfrac{f(\text{x})\,\text{dx}}{(1+\text{x}^2)}=\dfrac{\text{m}\pi}{\text{n}}$ (where m and n are relative prime) then value of (m + n) is
\qn{2} Let f(x) = 2x$^3$ $-$ 3x$^2$ $-$ x + $\dfrac{3}{2}$ and $\displaystyle\int_{1/8}^{7/8} f(f(\text{x}))\,\text{dx}=\dfrac{\text{m}}{\text{n}}$ (where m, n are relative prime) then value of (m + n) is equal to
\qn{3} f and h are two real valued continuous function and let $\displaystyle\int \text{h}(\text{x})\,\text{dx}=f^{-1}(\text{x})$ and f(x) = x$^3$ + x + sin 2$\pi$x + 2 and the value of $\displaystyle\int_2^4 \text{x}\,\text{h}(\text{x})\,\text{dx}$ is `k', then find value of 4k.
\qn{4} A six faced die is biased such that 3 $\times$ P(a prime number) = 6 $\times$ P(a composite number)= 2 $\times$ P(1). Let X be a random variable that counts the number of times one gets a perfect square on some throws of this die. If the die is thrown eleven times, then the mean of X is:
\qn{5} Let the sequence is 4, 7, 1, 8, 9, 7, 6, ..... for n $>$ 2. The n$^\text{th}$ term of the sequence is units digit of the sum of the two previous terms. Let S$_\text{n}$ denote the sum of the first n terms of this sequence. The smallest value of n for which S$_\text{n}$ $>$ 10004 is k, then (k $-$ 1990) is,
\qn{6} The value of $\dfrac{\displaystyle\sum_{r=0}^{k}{}^\text{n}\text{C}_\text{r}\,{}^{\text{n}+\text{k}-\text{r}-1}\text{C}_{\text{n}-1}}{\displaystyle\sum_{r=1}^{k}2^\text{r}\,{}^\text{n}\text{C}_\text{r}\,{}^{\text{k}-1}\text{C}_{\text{r}-1}}$ (where n $\geq$ k) is equal to
\raggedcolumns
\papersec{2}
\partsec{1}{PHYSICS}
\sect{SECTION-I (i)}
\qn{1} Four non-conducting wires each with mass m = 1 kg and length $\ell$ = 2.0 m are joined to form a square. They have charge with charge per unit length 1 c/m and $-$1 c/m shown in figure. Initially the square is kept on x$-$y plane and there exists uniform electric field E = 4$\dfrac{\text{N}}{\text{C}}$ along Z-direction. If the square is released from rest :-\par
\qfig{p2_phy_s1_q01_1}
\begin{qoptions}
\opt{A}{It will perform oscillatory motion.}
\opt{B}{Its maximum angular speed will be $\sqrt{12\sqrt{2}}$ rad/sec.}
\opt{C}{Its maximum angular speed will be $\sqrt{8\sqrt{3}}$ rad/sec.}
\opt{D}{It will perform SHM.}
\end{qoptions}
\qn{2} A student is measuring resistance of a galvanometer uses a high resistance box (HRB) in series with the galvanometer and a low resistance box (LRB) in parallel to the galvanometer. He introduces a resistance R = 500 $\Omega$ from HRB for deflection $\theta$. When he closes the key K$_2$ also in the figure shown below and introduce a resistance R = 10 $\Omega$ from LRB, without changing resistance from HRB deflection of galvanometer changes to $\dfrac{3}{28}\theta$. Resistance of galvanometer can be :-\par
\qfig{p2_phy_s1_q02_1}
\begin{qoptions}
\opt{A}{200 $\Omega$}
\opt{B}{20 $\Omega$}
\opt{C}{50 $\Omega$}
\opt{D}{100 $\Omega$}
\end{qoptions}
\qn{3} A spherical ball of mass M moving with initial velocity V collide elastically with another ball of mass M which is L shaped rigid massless frame as shown in figure. The L shaped frame contains another mass M connected at other end.\par
\qfig{p2_phy_s1_q03_1}
\begin{qoptions}
\opt{A}{The speed of striking mass after collision $\dfrac{\text{V}}{7}$ backward.}
\opt{B}{The speed of striking mass after collision $\dfrac{\text{V}}{3}$ in the same direction.}
\opt{C}{After a time interval of $\dfrac{7\pi\text{L}}{4\text{V}}$ the frame comes to orientation shown in fig. (2).}
\opt{D}{After a time interval of $\dfrac{7\pi\text{L}}{8\text{V}}$ the frame comes to orientation shown in fig. (2).}
\end{qoptions}
\qn{4} There is an insect inside a cabin eyeing towards a thick glass plate P$_1$. Insect sees the images of light source across glass plate P$_1$ outside the cabin. Cabin is made of thick glass plate of refractive index $\mu$ = $\dfrac{3}{2}$ and thickness 3 cm. Insect is eyeing from the middle of the cabin and there is an identical glass plate P$_2$ on the other side as shown in the figure. (glass plates are partially reflective and consider only paraxial rays)\par
\qfig{p2_phy_s1_q04_1}
\begin{qoptions}
\opt{A}{The first image is seen at a distance of 14 cm from eye of insect.}
\opt{B}{The second image is seen at a distance of 18 cm from eye of insect.}
\opt{C}{Number of images seen by insect is 2.}
\opt{D}{Number of images seen by insect is 4.}
\end{qoptions}
\qn{5} A particle of mass m = 2 kg is connected from the origin `O' of an inertial frame with the help of an elastic card of force constant K = $\dfrac{2\text{N}}{\text{m}}$ and negligible relaxed length. In addition to the elastic force of the cord, a constant force $\vec{\text{F}}$ = 10$\hat{\text{i}}$N also acts on the particle. The particle is pulled to an arbitrary position (4, 3)m as shown in figure and released.\par
\qfig{p2_phy_s1_q05_1}
\begin{qoptions}
\opt{A}{particle oscillate between (4, 3)m and (6, $-$3)m.}
\opt{B}{particle oscillate between (4, 3)m and (8, $-$3)m.}
\opt{C}{time period of oscillations is 2$\pi$ sec.}
\opt{D}{motion is simple harmonic.}
\end{qoptions}
\qn{6} A long uniform solid metallic cylinder has cross sectional radius r. Surface of the cylinder is painted by black body material. Resistivity of the material of inner cylinder is $\rho$. Another long hollow coaxial cylinder of black body material has cross sectional radius R = 8r. Electric current I continuously passes through the axis of inner cylinder. Whole arrangement lies in vacuum. When steady state is reached temperature of outer shell is 500 K. Choose the correct answer.\par
\qfig{p2_phy_s1_q06_1}
\begin{qoptions}
\opt{A}{Temperature of surface of inner cylinder is 500$\sqrt{3}$ K.}
\opt{B}{Current through the inner cylinder is $\sqrt{\dfrac{\sigma}{\rho}16\pi^2\text{r}^3(500)^4}$.}
\opt{C}{Intensity of emitted radiation at distance 16r from axis is $\dfrac{\sigma}{2}(500)^4$.}
\opt{D}{None of these.}
\end{qoptions}
\qn{7} Consider the situation shown in figure. Source of sound `S' and listener `O' both are moving towards a wall with velocities V$_\text{s}$ = 20 m/s and V$_0$ = 10 m/s respectively. The wall is moving away from them with velocity u = 5 m/s. The wind is blowing with velocity v$_\text{w}$ = 30 m/s towards the wall. The frequency of sound emitted by source when it is stationary is f$_0$ = 500 Hz. If velocity of sound relative to air is v = 300 m/s then :-\par
\qfig{p2_phy_s1_q07_1}
\begin{qoptions}
\opt{A}{frequency of sound received by listener after reflection from wall is 534 Hz approx.}
\opt{B}{wavelength of sound received by wall is 0.56 m.}
\opt{C}{wavelength of sound received by listener after reflection from wall is 0.57 m approx.}
\opt{D}{frequency of sound received by listener after reflection from wall is 524 Hz approx.}
\end{qoptions}
\qn{8} A surface consisting of curve inclined and two semicircular parts in vertical plane is as shown in diagram. There is no friction anywhere and ball of mass m smoothly crosses from one semicircular part to other at C. [Radius of ball is negligible.]\par
\qfig{p2_phy_s1_q08_1}
\begin{qoptions}
\opt{A}{The minimum value of h so that ball leaves the surface from point D is $\dfrac{5}{2}$R.}
\opt{B}{If ball just leaves at D the normal reaction by surface just after crossing C is 17 mg.}
\opt{C}{After leaving the surface at D at the distance from B ball strike the horizontal surface is $\left(\sqrt{\dfrac{51}{16}}\right)$R.}
\opt{D}{None of these.}
\end{qoptions}
\sect{SECTION-I (ii)}
\paraq{9}{10}Two objects of volume V$_1$ and V$_2$ and densities $\rho_1$ and $\rho_2$ have identical flat portion of area `S' on their surfaces. These flat portions are glued to each other. The composite body thus formed, floats fully submerged in a liquid with the common flat portion horizontal as shown in figure. The glue can with stand a maximum force F$_\text{max}$. If h = minimum depth in the liquid can the common flat portion be keeping the object intact in equilibrium and acceleration of free fall g = 10 m/s$^2$, list-I describe four physical quantities which have to be calculated and match it with list-II.\par
{\footnotesize\setlength{\tabcolsep}{2pt}\noindent\begin{tabular}{|l|L{4.2cm}|l|l|}
\hline
 & \multicolumn{1}{c|}{List - I} & & List-II\\ \hline
(I) & Value of h & (P) & 10 m\\ \hline
(II) & Force by liquid on upper object & (Q) & 5500 N\\ \hline
(III) & Force by liquid on lower object & (R) & 4500 N\\ \hline
(IV) & Total force by liquid on both object & (S) & 10$^4$ N\\ \hline
 & & (T) & 3000 N\\ \hline
 & & (U) & 7500 N\\ \hline
\end{tabular}}\par
\qn{9} If objects shown in figure has volume V$_1$ = 1 m$^3$ and V$_2$ = 1 m$^3$ densities $\rho_1$ = 400 kg/m$^3$ and $\rho_2$ = 600 kg/m$^3$ and identical flat portion has area S = 100 cm$^2$ and if the glue can withstand a maximum force F$_\text{max}$ = 500 N, then :- (Neglect atmospheric pressure).\par
\qfig{p2_phy_s2_q09_1}
\begin{qoptions}
\opt{A}{I$-$P, II$-$T, III$-$U, IV$-$S}
\opt{B}{I$-$P, II$-$R, III$-$T, IV$-$S}
\opt{C}{I$-$P, II$-$R, III$-$Q, IV$-$S}
\opt{D}{I$-$P, II$-$Q, III$-$R, IV$-$S}
\end{qoptions}
\qn{10} If object shown in figure has volume V$_1$ = 1m$^3$ and V$_2$ = 0.5 m$^3$ densities $\rho_1$ = 400 kg/m$^3$ and $\rho_2$ = 700 kg/m$^3$ and identical flat portion has area S = 100 cm$^2$, and if the glue can with stand a maximum force F$_\text{max}$ = 500 N (Neglect atmospheric pressure).\par
\qfig{p2_phy_s2_q10_1}
\begin{qoptions}
\opt{A}{I$-$P, II$-$T, III$-$R, IV$-$U}
\opt{B}{I$-$Q, II$-$T, III$-$R, IV$-$U}
\opt{C}{I$-$P, II$-$R, III$-$Q, IV$-$U}
\opt{D}{I$-$P, II$-$R, III$-$T, IV$-$U}
\end{qoptions}
\paraq{11}{12}A parallel plate capacitor has two sections. In these sections areas of plates are A$_1$ and A$_2$ and distances between the plates are d$_1$ and d$_2$ respectively. The distances between the plates are much smaller than their linear dimensions. An uniform electric field E that is perpendicular to the plates is switched on. Neglecting edge effects, E$_1$ and E$_2$ are electric fields established in sections (1) and (2) respectively.\par
\qfig{p2_phy_s2_para11_12_1}
List-I describe various physical quantities (in SI unit) match it with List-II.\par
{\footnotesize\setlength{\tabcolsep}{2pt}\noindent\begin{tabular}{|l|L{3.2cm}|l|l|}
\hline
 & \multicolumn{1}{c|}{List-I} & & List-II\\ \hline
(I) & Value of E$_1$ & (P) & 150\\ \hline
(II) & Value of E$_2$ & (Q) & 75\\ \hline
(III) & Value of E$_1$ + E$_2$ & (R) & 225\\ \hline
(IV) & Value of E$_1$ $-$ E$_2$ & (S) & 175\\ \hline
 & & (T) & 450\\ \hline
 & & (U) & 300\\ \hline
\end{tabular}}\par
\qn{11} If values A$_1$ = 1 m$^2$ and A$_2$ = 2 m$^2$, d$_1$ = 1 mm and d$_2$ = 2 mm and applied field is 100 N/c match List I with List II.
\begin{qoptions}
\opt{A}{I$-$P, II$-$Q, III$-$R, IV$-$Q}
\opt{B}{I$-$R, II$-$Q, III$-$U, IV$-$T}
\opt{C}{I$-$R, II$-$S, III$-$U, IV$-$T}
\opt{D}{I$-$P, II$-$Q, III$-$R, IV$-$T}
\end{qoptions}
\qn{12} If A$_1$ = 4m$^2$, A$_2$ = 8 m$^2$, d$_1$ = 6 mm an d$_2$ = 12 mm and applied field is 100 N/c. Match List I with List II.
\begin{qoptions}
\opt{A}{I$-$R, II$-$Q, III$-$U, IV$-$T}
\opt{B}{I$-$P, II$-$Q, III$-$R, IV$-$Q}
\opt{C}{I$-$P, II$-$Q, III$-$R, IV$-$T}
\opt{D}{I$-$R, II$-$S, III$-$U, IV$-$T}
\end{qoptions}
\sect{SECTION-II}
\qn{1} The figure shows an uniform slab of mass 7kg. Find moment of inertia (in kg cm$^2$) about the axis normal to the plane of drawing and passing through the point O. (Given OA = OC, BC = 10cm and AB = 2 cm.)\par
\qfig{p2_phy_s3_q01_1}
\qn{2} Fixed thin wire frame in the form of a square with a side a is located in the horizontal plane in a uniform magnetic field, the induction of which is equal to B$_0$ and perpendicular to the plane of the frame. There is a wire jumper on the frame PP$_1$ having a mass m as shown in the figure. The frame and jumper are made from the same resistive wire $\rho$ and cross-sectional area S. The speed acquired by the jumper immediately after the magnetic field is turned off is given by v = $\dfrac{\text{a}^2\text{SB}_0^2}{\alpha\text{m}\rho}$. The value of $\alpha$ is (Neglect the friction force and the displacement of the jumper during the decay of the field.)\par
\qfig{p2_phy_s3_q02_1}
\qn{3} The network shown consists of an infinitely large number of identical sections. Resistance of each branch is R = $\sqrt{3}\,\Omega$. If equivalent resistance of the whole network between A and B is n$\Omega$, find n.\par
\qfig{p2_phy_s3_q03_1}
\qn{4} One mole of an ideal mono-atomic gas undergoes a process A $\to$ B that is a straight line on PV indicator diagram of the process as shown in the figure. Find volume of the gas in litre when the process turns from an endothermic to an exothermic one.\par
\qfig{p2_phy_s3_q04_1}
\qn{5} A uniform thin rod has mass m and length $\ell$. One end of the rod lies over rough horizontal surface and other end is connected to a vertical string, as shown in the figure. When string is cut, there is no slipping between rod and surface. Calculate the friction on the rod immediately after the string is cut. (given mg = $\dfrac{16}{\sqrt{3}}$)\par
\qfig{p2_phy_s3_q05_1}
\qn{6} A homogeneous thin rod of mass M and length L. At a height L from the hinge a thin weightless inextensible thread is attached vertically above the place where the rod is attached to the wall, the second end of which is attached to the free end of the rod (see figure). A thin weightless liquid film is stretched between the wall, the rod and the thread, the coefficient of surface tension of which is equal to $\sigma$. The rod is in a state of equilibrium and is located horizontally. Find the length of the thread $\ell$, assuming that it is less than $\pi$L/2. Take : $\sigma$ = $\dfrac{\text{Mg}}{9\sqrt{6}}$ and L = $\dfrac{9}{\sqrt{2}\pi}$ m.\par
\qfig{p2_phy_s3_q06_1}
\raggedcolumns
\partsec{2}{CHEMISTRY}
\sect{SECTION-I (i)}
\qn{1} For the given complexes the correct statement is/are\par
[Mo(CO)$_3$(PF$_3$)$_3$], [Mo(CO)$_3$(PCl$_3$)$_3$],\par [Mo(CO)$_3$(PMe$_3$)$_3$],
\begin{qoptions}
\opt{A}{[Mo(CO)$_3$(PF$_3$)$_3$] has max bond length of CO bond}
\opt{B}{[Mo(CO)$_3$(PMe$_3$)$_3$] has max bond length of CO bond}
\opt{C}{[Mo(CO)$_3$(PMe$_3$)$_3$] has minimum $\pi$-character in M-C bond}
\opt{D}{[Mo(CO)$_3$(PF$_3$)$_3$] has max bond order of CO bond}
\end{qoptions}
\qn{2} Which of the following is/are CORRECT statements?
\begin{qoptions}
\opt{A}{Ore bauxite is purified by alkaline treatment in the Baeyer process and then dissolved in molten cryolite and reduced electrolytically in the fused salt system to get metal}
\opt{B}{Sulphide ore of some of the less electropositive metals like Hg, Cu, Pb etc. are heated in air, a part of these is changed into oxide or sulphate then that reacts with the remaining part of the sulphide ore to give its metal and SO$_2$}
\opt{C}{Metallic gold is dissolved from its ore in dilute NaCN solution in presence of air and the solution so obtained is treated with scrap zinc when gold metal is precipitated}
\opt{D}{Magnesium is precipitated as Mg(OH)$_2$ from sea water by treatment with lime and then it dissolved in dilute HCl and reduced electrolytically in aqueous medium to get magnesium metal}
\end{qoptions}
\qn{3} Which of the following will give negative Kjeldahl and Lassaigne test for Nitrogen.
\begin{qoptions}
\opt{A}{Aniline}
\opt{B}{Hydrazine}
\opt{C}{Benzenediazonium chloride}
\opt{D}{Dimethyl ethanamine}
\end{qoptions}
\qn{4} Incorrect option among the following is/are
\begin{qoptions}
\opt{A}{Nucleophilicity order $\Rightarrow$ PhS$^\ominus$ $>$ CH$_3^\ominus$ $>$ I$^\ominus$ $>$ EtO$^\ominus$}
\opt{B}{Acidic strength order $\Rightarrow$ CH$_3$CO$_2$H $>$ H$_2$CO$_3$ $>$ \qfigi{p2_chem_s1_q04_1}}
\opt{C}{Basic strength order $\Rightarrow$ \qfigi{p2_chem_s1_q04_2}}
\opt{D}{Stability order : $\overset{+}{\text{C}}\text{Cl}_3$ $>$ $\overset{+}{\text{C}}\text{F}_3$ $>$ $\overset{+}{\text{C}}\text{H}_3$ $>$ CMe$_3$}
\end{qoptions}
\qn{5} d$_{\text{H}_2\text{O}}$ = 1 g/ml; \quad d$_{\text{C}_2\text{H}_5\text{OH}}$ = 0.92 g/ml\par
540 ml of H$_2$O is mixed with 500 ml of C$_2$H$_5$OH to form an ideal solution. Consider solute present to be non-volatile and non-electrolyte (K$_\text{b}$ (H$_2$O) = 0.54 K kg mole$^{-1}$; K$_\text{b}$ (C$_2$H$_5$OH) = 1.2 K kg mole$^{-1}$)\par
Boiling point of water = 373 K;\par
Vapour pressure of H$_2$O = 35 mm Hg\par
Vapour pressure of C$_2$H$_5$OH = 40 mm Hg
\begin{qoptions}
\opt{A}{Boiling point of above solution is = 383 K}
\opt{B}{If 500 ml C$_2$H$_5$OH is further added to above solution then the boiling point elevation of new solution is 36 K}
\opt{C}{Initial vapour pressure of solution 26.25 mm Hg}
\opt{D}{Initial vapour pressure of solution 8.75 mm Hg}
\end{qoptions}
\qn{6} Which of the following statement is/are CORRECT?
\begin{qoptions}
\opt{A}{Higher is the gold number of a lyophilic sol. more is its coagulating power}
\opt{B}{An electrolyte can coagulate both positive sol and negative sol}
\opt{C}{Impurities present in a colloid makes it more stable}
\opt{D}{Dialysis is a process to remove impurities of ions and smaller molecules from a colloid}
\end{qoptions}
\qn{7} The nucleus of an atom is located at x = y = z = 0. If the probability of finding an `s' orbital electron in a tiny volume around x = a, y = z = 0 is 1.0 $\times$ 10$^{-5}$. Choose the correct statement(s) regarding the probability of finding electron.
\begin{qoptions}
\opt{A}{The probability of finding the electron in the same-sized volume around x = z = 0, y = a is the same i.e., 1 $\times$ 10$^{-5}$}
\opt{B}{The probability of finding the electron in the same-sized volume around x = y = 0, z = a is zero}
\opt{C}{The probability of finding the electron at x = y = z = 0 is zero.}
\opt{D}{The probability of finding the electron at x = 2a, y = z = 0 must be greater than 1.0 $\times$ 10$^{-5}$.}
\end{qoptions}
\qn{8} Ca$_3$P$_2$ + H$_2$O $\longrightarrow$ A$\uparrow$ + B\par
Which of the following statements is/are correct regarding A and B?
\begin{qoptions}
\opt{A}{Gas `A' is passed into CuSO$_4$ solution and black ppt. is obtained}
\opt{B}{CO$_2$ is passed into excess of `B' and clear solution is obtained}
\opt{C}{When ``A'' is pure burns in air spontaneously}
\opt{D}{`A' is slightly soluble in water}
\end{qoptions}
\sect{SECTION-I (ii)}
\paraq{9}{10}Answer the following by appropriately matching the lists based on the information given in the paragraph.\par
\qfig{p2_chem_s2_para9_10_1}
\qn{9} Incorrect Match:
\begin{qoptions}
\opt{A}{(I) $\longrightarrow$ (T)}
\opt{B}{(II) $\longrightarrow$ (R)}
\opt{C}{(III) $\longrightarrow$ (S)}
\opt{D}{(IV) $\longrightarrow$ (P)}
\end{qoptions}
\qn{10} Correct Match:
\begin{qoptions}
\opt{A}{(I) $\longrightarrow$ (R)}
\opt{B}{(II) $\longrightarrow$ (Q)}
\opt{C}{(III) $\longrightarrow$ (S)}
\opt{D}{(IV) $\longrightarrow$ (P)}
\end{qoptions}
\paraq{11}{12}Answer the following by appropriately matching the lists based on the information given in the paragraph.\par
{\footnotesize\setlength{\tabcolsep}{1.5pt}\noindent\begin{tabular}{|l|L{2.2cm}|l|L{3.4cm}|}
\hline
 & \multicolumn{1}{c|}{List - I} & & \multicolumn{1}{c|}{List - II}\\ \hline
(I) & K$_4$[Fe(CN)$_6$] & (P) & Produces deep blue solution with CuSO$_4$ solution\\ \hline
(II) & (NH$_4$)$_2$S & (Q) & Produces deep blue PPT with FeSO$_4$ solution\\ \hline
(III) & NH$_4$OH (excess) & (R) & Produces deep blue PPT with FeCl$_3$ solution\\ \hline
(IV) & NaOH (excess) & (S) & Produces colourless solution or white PPT with ZnSO$_4$ solution\\ \hline
 & & (T) & Produces chocolate brown PPT with CuSO$_4$ solution\\ \hline
 & & (U) & Produces reddish brown precipitate with FeCl$_3$\\ \hline
\end{tabular}}\par
\qn{11} Which of the following is incorrect match?
\begin{qoptions}
\opt{A}{I $\longrightarrow$ R}
\opt{B}{I $\longrightarrow$ S}
\opt{C}{I $\longrightarrow$ Q}
\opt{D}{I $\longrightarrow$ T}
\end{qoptions}
\qn{12} Which of the following is correct match?
\begin{qoptions}
\opt{A}{IV $\longrightarrow$ S,U}
\opt{B}{III $\longrightarrow$ P,S,U}
\opt{C}{II $\longrightarrow$ S}
\opt{D}{All of the above}
\end{qoptions}
\sect{SECTION-II}
\qn{1} Aqueous urea solution is hydrolysed according to the reaction\par
$(\text{NH}_2)_2\text{CO}_{(\text{aq.})} + \text{H}_2\text{O}_{(\ell)} \rightarrow 2\text{NH}_{3(\text{aq.})} + \text{CO}_{2(\text{aq.})}$\par
the standard enthalpy of formation and entropies at 300 K for reactants and products are\par
{\footnotesize\setlength{\tabcolsep}{2pt}\noindent\begin{tabular}{|l|c|c|c|c|}
\hline
 & CO$_{2(\text{aq.})}$ & H$_2$O$_{(\ell)}$ & NH$_{3(\text{aq.})}$ & (NH$_2$)$_2$CO$_{(\text{aq.})}$\\ \hline
$\Delta\text{H}^0$ (kJ/mole) & $-$412.9 & $-$285.8 & $-$80.8 & $-$317.7\\ \hline
S$^\circ$(JK$^{-1}$mole$^{-1}$) & 121.0 & 69.9 & 110.0 & 176.0\\ \hline
\end{tabular}}\par
Find value of lnK of above reaction at 300 K.
\qn{2} 2.24 g of a mixture having impure Fe$_2$(SO$_4$)$_3$ sample and pure FeC$_2$O$_4$ was dissolved in dil. H$_2$SO$_4$. the complete oxidation of reaction mixture requires 100 ml $\dfrac{\text{N}}{20}$KMnO$_4$ solution. After the oxidation the reaction mixture was reduced completely into Fe$^{+2}$ by Zn and dil. H$_2$SO$_4$. This reduced solution requires 200 ml of same KMnO$_4$ solution. Then find percentage purity of Fe$_2$(SO$_4$)$_3$ sample. [Fe = 56]
\qn{3} The decomposition of A$_2$B$_2$ follows the rate equation r = K[A$_2$B$_2$]$^2$.\par
A$_2$B$_2$ $\xrightarrow{\text{K}}$ A$_2$ + B$_2$. The kinetic data of the reaction is given below.\par
{\footnotesize\setlength{\tabcolsep}{2pt}\noindent\begin{tabular}{|L{5cm}|c|}
\hline
Initial conc. of A$_2$B$_2$ (mole/lit) & 7\\ \hline
First half life of A$_2$B$_2$ in sec. & 1440\\ \hline
\end{tabular}}\par
If the initial concentration of A$_2$B$_2$ were 56 m mol/lit, Calculate the time (in sec) taken for the conversion of A$_2$B$_2$ to reach 60\% of initial value.\par
(Report your answer as $\dfrac{\text{t}}{1000}$ )
\qn{4} Reagents which can convert CH$_3$$-$CHO into carboxylic acid/salt :\par
(i) [Ag(NH$_3$)$_2$]OH \quad (ii) CH$_3$$-$CO$_3$H \quad (iii) NaOI \quad (iv) DIBAL$-$H\par
(v) Cu/300\textdegree C \quad (vi) K$_2$Cr$_2$O$_7$/H$^+$ \quad (vii) Al$\left(-\text{O}-\text{CH}(\text{CH}_3)_2\right)_3$ / $\text{H}_3\text{C}-\underset{\text{CH}_3}{\text{CH}}-\text{OH}$
\qn{5} Calculate the isoelectric point (pI) of lysine. Give the pK$_\text{a}$ of $\alpha-\overset{\oplus}{\text{N}}\text{H}_3$ is 8.95, pK$_\text{a}$ side chain $\overset{\oplus}{\text{N}}\text{H}_3$ is 10.53 and pK$_\text{a}$ of $\alpha$-COOH is 2.18
\qn{6} Find sum of the number of B$-$O$-$B bonds in\par
(i) [B$_4$O$_5$(OH)$_4$]$^{2-}$\par
(ii) [B$_5$O$_6$(OH)$_4$]$^{-}$
\raggedcolumns
\partsec{3}{MATHEMATICS}
\sect{SECTION-I (i)}
\qn{1} There are three distinct points A(2, 0); B and C such that OB = OC = 2. Consider the triangle T formed by the points A, B and C. Let there are exactly m such triangles T that are equilateral and there are exactly n such triangles T that are right angled isosceles. Then 2m + n is
\begin{qoptions}
\opt{A}{divisible by 2}
\opt{B}{Composite number}
\opt{C}{prime number}
\opt{D}{divisible by 5}
\end{qoptions}
\qn{2} If f(x) = 3$^{-|\text{x}|}$; x $\in$ R and h(t) = $\displaystyle\int_{-1}^{1} f\left(\dfrac{\text{x}}{\text{t}}\right)\text{dx}$; t $\neq$ 0 then which of the following is/are true
\begin{qoptions}
\opt{A}{$\displaystyle\lim_{\text{t}\to0^+}\dfrac{\text{h(t)}}{\text{t}}=\dfrac{2}{\ln 3}$}
\opt{B}{$\displaystyle\lim_{\text{t}\to0^-}\dfrac{\text{h(t)}}{\text{t}}=\dfrac{-2}{\ln 3}$}
\opt{C}{$\displaystyle\lim_{\text{t}\to0}\dfrac{\text{h(t)}}{\text{t}}$ does not exist}
\opt{D}{$\displaystyle\lim_{\text{t}\to0}\dfrac{\text{h(t)}}{\text{t}}=\dfrac{2}{\ln 3}$}
\end{qoptions}
\qn{3} Let g : R $\to$ R be a continuous function such that g(x + 1) = $\dfrac{1}{3}$g(x) for all x $\in$ R and let $\alpha_\text{n}$ = $\displaystyle\int_0^\text{n}$ g(x)dx for all integer n $\geqslant$ 1. Then $\displaystyle\lim_{\text{n}\to\infty}\alpha_\text{n}$
\begin{qoptions}
\opt{A}{Exist}
\opt{B}{Does not exist}
\opt{C}{is equal to $\dfrac{3}{2}\displaystyle\int_0^1$ g(x)dx}
\opt{D}{not equal to $\dfrac{3}{2}\displaystyle\int_0^1$ g(x)dx}
\end{qoptions}
\qn{4} t$_1$ = $\left(\sin^{-1}\text{x}\right)^{\sin^{-1}\text{x}}$; t$_2$ = $\left(\sin^{-1}\text{x}\right)^{\cos^{-1}\text{x}}$; t$_3$ = $\left(\cos^{-1}\text{x}\right)^{\sin^{-1}\text{x}}$; t$_4$ = $\left(\cos^{-1}\text{x}\right)^{\cos^{-1}\text{x}}$. Where sgn(x) represents signum function. Then which of the following statements is/are correct.
\begin{qoptions}
\opt{A}{If x $\in$ (0, cos1) then sgn(t$_4$ $-$ t$_2$ + t$_3$ $-$ t$_1$) = 1}
\opt{B}{If x $\in$ $\left(\cos 1,\dfrac{1}{\sqrt{2}}\right)$ then log$_{10}$(t$_3$ $-$ t$_1$ + t$_4$ $-$ t$_2$) is defined}
\opt{C}{x $\in$ (sin1, 1) then t$_1$ $>$ t$_2$ $>$ t$_4$ $>$ t$_3$}
\opt{D}{x $\in$ $\left(\dfrac{1}{\sqrt{2}},\sin 1\right)$ then t$_2$ $>$ t$_1$ $>$ t$_3$ $>$ t$_4$}
\end{qoptions}
\qn{5} M is a 3 $\times$ 3 matrix with integers entries. For M we have sum of entries in column 2 is equal to 4 times the sum of entries in column 1, sum of entries in column 3 is 4 times the sum of entries in column 2. Also sum of entries in Row 2 exceeds the sum of entries in row 1 by 6 and sum of entries in row 3 exceeds the sum of entries in row 2 by 6. Then which of the following is/are incorrect.
\begin{qoptions}
\opt{A}{The sum of all the entries in M must be divisible by 3 only}
\opt{B}{None of the row sums is divisible by 7}
\opt{C}{One of the column sums must be divisible by 7}
\opt{D}{The sum of all the entries in M must be divisible by 21}
\end{qoptions}
\qn{6} Consider the equation Z$^{200}$ $-$ $\lambda$Z$^{199}$ + 1 = 0. Let Z$_1$, Z$_2$, Z$_3$ ..... Z$_{200}$ be non-zero complex numbers and Z$^{200}$ $-$ $\lambda$Z$^{199}$ + 1 = (Z $-$ Z$_1$) (Z $-$ Z$_2$) (Z $-$ Z$_3$) ..... (Z $-$ Z$_{200}$) then
\begin{qoptions}
\opt{A}{Imaginary part of $\displaystyle\sum_{j=1}^{200}\sum_{k=1}^{200}\dfrac{\text{Z}_\text{j}}{\text{Z}_\text{k}}$ is 0}
\opt{B}{Imaginary part of $\displaystyle\sum_{j=1}^{200}\sum_{k=1}^{200}\dfrac{\text{Z}_\text{j}}{\text{Z}_\text{k}}$ is 1}
\opt{C}{Real part of $\displaystyle\sum_{j=1}^{200}\sum_{k=1}^{200}\dfrac{\text{Z}_\text{j}}{\text{Z}_\text{k}}$ is 0}
\opt{D}{Real part of $\displaystyle\sum_{j=1}^{200}\sum_{k=1}^{200}\dfrac{\text{Z}_\text{j}}{\text{Z}_\text{k}}$ is 1}
\end{qoptions}
\qn{7} Consider the determinant\par\noindent\resizebox{\linewidth}{!}{$\text{D}=\begin{vmatrix}\dfrac{1-\ell_1^3\text{m}_1^3}{1-\ell_1\text{m}_1}&\dfrac{1-\ell_2^3\text{m}_1^3}{1-\ell_2\text{m}_1}&\dfrac{1-\ell_3^3\text{m}_1^3}{1-\ell_3\text{m}_1}\\[2ex]\dfrac{1-\ell_1^3\text{m}_2^3}{1-\ell_1\text{m}_2}&\dfrac{1-\ell_2^3\text{m}_2^3}{1-\ell_2\text{m}_2}&\dfrac{1-\ell_3^3\text{m}_2^3}{1-\ell_3\text{m}_2}\\[2ex]\dfrac{1-\ell_1^3\text{m}_3^3}{1-\ell_1\text{m}_3}&\dfrac{1-\ell_2^3\text{m}_3^3}{1-\ell_2\text{m}_3}&\dfrac{1-\ell_3^3\text{m}_3^3}{1-\ell_3\text{m}_3}\end{vmatrix}$}\par
here $\ell_1$, $\ell_2$, $\ell_3$, m$_1$, m$_2$, m$_3$ $\in$ R and $\ell_\text{i}$m$_\text{j}$ $\neq$ 1 $\forall$ i, j $\in$ \{1, 2, 3\} also $\ell_1<\ell_2<\ell_3$ \& m$_1<$ m$_2<$ m$_3$, then choose the INCORRECT option(s).
\begin{qoptions}
\opt{A}{D $>$ 0}
\opt{B}{D $<$ 0}
\opt{C}{D = 0}
\opt{D}{D = $-$1}
\end{qoptions}
\qn{8} Consider a parallelogram ABCD. A straight line is drawn which cuts AB, AC, AD at points P, Q, R respectively. If $\overrightarrow{\text{AP}}$ = $\alpha\overrightarrow{\text{AB}}$, $\overrightarrow{\text{AR}}$ = $\beta\overrightarrow{\text{AD}}$ and $\overrightarrow{\text{AQ}}$ = $\dfrac{\gamma}{2}\overrightarrow{\text{AC}}$ where $\alpha$, $\beta$, $\gamma$ are positive scalars, then
\begin{qoptions}
\opt{A}{$\alpha$, $\beta$, $\gamma$ are in A.P.}
\opt{B}{$\alpha$, $\beta$, $\gamma$ are in H.P.}
\opt{C}{$\alpha$, $\gamma$, $\beta$ are in H.P.}
\opt{D}{$\alpha\beta\geq\gamma^2$}
\end{qoptions}
\sect{SECTION-I (ii)}
\qn{9} 3 circles C$_1$, C$_2$ and C$_3$ with centres O$_1$, O$_2$ and O$_3$ of radius 9, 6 and 3 respectively touches each other as shown in figure. A chord AB to circle C$_1$ touches circles C$_2$ and C$_3$ at P and Q respectively. Line PQ meets line joining centres of C$_2$ and C$_3$ at T, D is mid-point of Q and O$_2$.\par
\qfig{p2_math_s2_q09_1}
{\footnotesize\setlength{\tabcolsep}{2pt}\noindent\begin{tabular}{|l|L{3.5cm}|l|l|}
\hline
 & \multicolumn{1}{c|}{List - I} & & List - II\\ \hline
(I) & Length of chord PQ is & (P) & $3\sqrt{2}$\\ \hline
(II) & Length of chord AB is & (Q) & $4\sqrt{14}$\\ \hline
(III) & Length of TP is & (R) & $6\sqrt{2}$\\ \hline
(IV) & Length of O$_3$D is & (S) & $4\sqrt{7}$\\ \hline
 & & (T) & $2\sqrt{14}$\\ \hline
 & & (U) & $12\sqrt{2}$\\ \hline
\end{tabular}}\par
Which of the following is the only correct combination?
\begin{qoptions}
\opt{A}{I $\longrightarrow$ P}
\opt{B}{I $\longrightarrow$ U}
\opt{C}{II $\longrightarrow$ Q}
\opt{D}{II $\longrightarrow$ S}
\end{qoptions}
\qn{10} 3 circles C$_1$, C$_2$ and C$_3$ with centres O$_1$, O$_2$ and O$_3$ of radius 9, 6 and 3 respectively touches each other as shown in figure. A chord AB to circle C$_1$ touches circles C$_2$ and C$_3$ at P and Q respectively. Line PQ meets line joining centres of C$_2$ and C$_3$ at T, D is mid-point of Q and O$_2$.\par
\qfig{p2_math_s2_q10_1}
{\footnotesize\setlength{\tabcolsep}{2pt}\noindent\begin{tabular}{|l|L{3.5cm}|l|l|}
\hline
 & \multicolumn{1}{c|}{List - I} & & List - II\\ \hline
(I) & Length of chord PQ is & (P) & $3\sqrt{2}$\\ \hline
(II) & Length of chord AB is & (Q) & $4\sqrt{14}$\\ \hline
(III) & Length of TP is & (R) & $6\sqrt{2}$\\ \hline
(IV) & Length of O$_3$D is & (S) & $4\sqrt{7}$\\ \hline
 & & (T) & $2\sqrt{14}$\\ \hline
 & & (U) & $12\sqrt{2}$\\ \hline
\end{tabular}}\par
Which of the following is the only incorrect combination?
\begin{qoptions}
\opt{A}{II $\longrightarrow$ Q}
\opt{B}{III $\longrightarrow$ U}
\opt{C}{IV $\longrightarrow$ P}
\opt{D}{IV $\longrightarrow$ T}
\end{qoptions}
\qn{11} Let $f$(x) = $\begin{cases}[2\text{x}]; & -1\leq\text{x}<1\\ |2\text{x}|; & 1\leq\text{x}\leq2\end{cases}$, g(x) = $\dfrac{\text{cosecx}}{2}$; x $\in$ (0, 2$\pi$) $-$ \{$\pi$\} (where [$\cdot$] denotes the greatest integer function).\par
List-I : Contains information about $f$($f$(x)), $f$(g(x)), $f$(2g(x)), g($f$(x)).\par
List-II : Contains information about the continuity and differentiability of $f$($f$(x)), $f$(g(x)), $f$(2g(x)), $f$(2g(x)), g($f$(x)).\par
{\footnotesize\setlength{\tabcolsep}{1.5pt}\noindent\begin{tabular}{|l|L{3.4cm}|l|L{2.7cm}|}
\hline
 & \multicolumn{1}{c|}{List - I} & & \multicolumn{1}{c|}{List - II}\\ \hline
(I) & $f$(g(x)) = $-$2 $\forall$ x $\in$ $\left[\dfrac{5\pi}{4},\dfrac{5\pi}{3}\right]-\left\{\dfrac{3\pi}{2}\right\}$ & (P) & discontinuous at x = 1\\ \hline
(II) & g($f$(x)) = 1 for some x $\in$ (1, 2) & (Q) & continuous at x = $\dfrac{\pi}{2}$\\ \hline
(III) & $f$($f$(x)) = 0 for $\forall$ x $\in$ $\left[\dfrac{1}{4},\dfrac{1}{3}\right]$ & (R) & discontinuous at x = $\dfrac{5\pi}{6}$\\ \hline
(IV) & $f$(2g(x)) = 2 for some x $\in$ $\left[\dfrac{\pi}{6},\dfrac{5\pi}{6}\right]$ & (S) & continuous at x = $\dfrac{\pi}{6}$\\ \hline
 & & (T) & differentiable at x = $\dfrac{\pi}{2}$\\ \hline
 & & (U) & differentiable at x = $\dfrac{\pi}{6}$\\ \hline
\end{tabular}}\par
Which of the following is the only CORRECT combination?
\begin{qoptions}
\opt{A}{I $\longrightarrow$ P,T}
\opt{B}{I $\longrightarrow$ R,T}
\opt{C}{II $\longrightarrow$ Q,S}
\opt{D}{II $\longrightarrow$ P,Q}
\end{qoptions}
\qn{12} Let $f$(x) = $\begin{cases}[2\text{x}]; & -1\leq\text{x}<1\\ |2\text{x}|; & 1\leq\text{x}\leq2\end{cases}$, g(x) = $\dfrac{\text{cosecx}}{2}$; x $\in$ (0, 2$\pi$) $-$ \{$\pi$\} (where [$\cdot$] denotes the greatest integer function).\par
List-I : Contains information about $f$($f$(x)), $f$(g(x)), $f$(2g(x)), g($f$(x)).\par
List-II : Contains information about the continuity and differentiability of $f$($f$(x)), $f$(g(x)), $f$(2g(x)), $f$(2g(x)), g($f$(x)).\par
{\footnotesize\setlength{\tabcolsep}{1.5pt}\noindent\begin{tabular}{|l|L{3.4cm}|l|L{2.7cm}|}
\hline
 & \multicolumn{1}{c|}{List - I} & & \multicolumn{1}{c|}{List - II}\\ \hline
(I) & $f$(g(x)) = $-$2 $\forall$ x $\in$ $\left[\dfrac{5\pi}{4},\dfrac{5\pi}{3}\right]-\left\{\dfrac{3\pi}{2}\right\}$ & (P) & discontinuous at x = 1\\ \hline
(II) & g($f$(x)) = 1 for some x $\in$ (1, 2) & (Q) & continuous at x = $\dfrac{\pi}{2}$\\ \hline
(III) & $f$($f$(x)) = 0 for $\forall$ x $\in$ $\left[\dfrac{1}{4},\dfrac{1}{3}\right]$ & (R) & discontinuous at x = $\dfrac{5\pi}{6}$\\ \hline
(IV) & $f$(2g(x)) = 2 for some x $\in$ $\left[\dfrac{\pi}{6},\dfrac{5\pi}{6}\right]$ & (S) & continuous at x = $\dfrac{\pi}{6}$\\ \hline
 & & (T) & differentiable at x = $\dfrac{\pi}{2}$\\ \hline
 & & (U) & differentiable at x = $\dfrac{\pi}{6}$\\ \hline
\end{tabular}}\par
Which of the following is the only CORRECT combination?
\begin{qoptions}
\opt{A}{IV $\longrightarrow$ R,T}
\opt{B}{IV $\longrightarrow$ P,S}
\opt{C}{III $\longrightarrow$ P,U}
\opt{D}{II $\longrightarrow$ R,T}
\end{qoptions}
\sect{SECTION-II}
\qn{1} The equation of plane of a $\Delta$ABC is $\dfrac{\text{x}}{\text{a}}+\dfrac{\text{y}}{\text{b}}+\dfrac{\text{z}}{\text{c}}=\lambda$ and lengths of sides of $\Delta$ABC opposite to vertices A, B, C are respectively p, q, r. Mid points of sides of $\Delta$ABC lies on coordinate axes. If a$^2$ = $\dfrac{\text{q}^2+\text{r}^2-\text{p}^2}{2}$, b$^2$ = $\dfrac{\text{p}^2+\text{r}^2-\text{q}^2}{2}$ and c$^2$ = $\dfrac{\text{p}^2+\text{q}^2-\text{r}^2}{2}$ then value of $\left|\dfrac{1}{\lambda}\right|$ is
\qn{2} Let S = \{1, 2, 3, 4\}. Then the number of elements in the set \{f : S $\times$ S $\longrightarrow$ S : f is onto and f(a, b) = f(b, a) $\geq$a $\forall$ (a, b)$\in$ S $\times$ S\} is
\qn{3} Let f be real function defined on R (the set of real numbers) such that\par
f$'$(x) = 100 (x $-$ 1) (x $-$ 2)$^2$(x $-$ 3)$^3$.... (x $-$ 100)$^{100}$, for all x $\in$ R. If g is a function defined on R such that $\displaystyle\int_\text{a}^\text{x}\text{e}^{f(\text{t})}\text{dt}=\int_0^\text{x}\text{g}(\text{x}-\text{t})\,\text{dt}+2\text{x}+3$, if sum of the all the values of x for which g(x) has a local extremum be $\lambda$ then find $\dfrac{\lambda}{500}$.
\qn{4} The number of pairs ordered (x, y) satisfying the equation sin x + sin y = sin (x + y) and $|$x$|$ + $|$y$|$ = 1 is equal to
\qn{5} Let A is non singular matrices of order 2 such that\par
9A$^2$B $-$ 6AB + B = 9$\left\{\text{B}-\begin{bmatrix}2&-3\\2&1\end{bmatrix}\right\}$. If B = $\begin{bmatrix}3&0\\1&1\end{bmatrix}$ and Trace (A) = $\dfrac{10}{3}$ then $|$A$|$ is
\qn{6} If the quadratic equation x$^2$ + (2 $-$ tan $\theta$) x $-$ (1 + tan $\theta$) = 0 has two integral roots and sum of all possible values of $\theta$ in the interval (0, 2$\pi$) is k$\pi$, then the value of k is equal to \_\_\_\_\_
\end{multicols}
\end{document}
